

\documentclass[5pt]{article}

\usepackage{sectsty}
\usepackage{graphicx}
\usepackage{lipsum} % for generating dummy text
\usepackage[margin=1in]{geometry}
\usepackage{setspace}
\usepackage{array}
\usepackage{cellspace}
\usepackage{tabularx}
\usepackage[table]{xcolor}
\usepackage{tabularray}
\usepackage{pgfplots}


\usepackage{hyperref}
\usepackage{scrextend}
\graphicspath{ {../assets} }

\usepackage[italian]{babel}


% package and setup for tables
\usepackage{float}
\usepackage[table]{xcolor}
\renewcommand{\arraystretch}{1.5}
\arrayrulecolor{black}


% Margins
\topmargin=-0.45in
\evensidemargin=0in
\oddsidemargin=0in
\textwidth=6.5in
\textheight=9.0in
\headsep=0.25in

\setlength{\parindent}{0pt}

\title{Piano di Qualifica}
\author{Jackpot Coding}
\renewcommand*\contentsname{Indice}
\date{\today}
%STARTOF THE DOCUMENT
\begin{document}
	
	%-------------------------
	
	% Reduce top margin only on the first page
	\newgeometry{top=0.5in}
	
	%UNIPD LOGO
	\vspace{8pt}
		\includegraphics[scale=0.2]{UNIPDFull.png}
	%END UNIPD LOGO
	
	\vspace{30pt}
	
	%COURSE INFO
	\begin{minipage}[t]{0.48\textwidth}
		%COURSE TITLE
		\begin{flushleft}
			Informatica\\
			\vspace{5pt}
			\textbf{\LARGE Ingegneria del Software}\\
			Anno Accademico: 2023/2024
		\end{flushleft}
		%END COURSE TITLE
	\end{minipage}
	%END COURSE INFO
	
	
	\vspace{5px}
	
	
	%BLACK LINE
	\rule{\textwidth}{5pt}
	
	%JACKPOT CODING INFO
	\begin{minipage}[t]{0.50\textwidth}
		%LOGO JACKPOT CODING
		\begin{flushleft}
			\hspace{10pt}
				\includegraphics[scale=0.65]{jackpot-logo.png} 
		\end{flushleft}
	\end{minipage}
	\hspace{-60pt} % This adds horizontal space between the minipages
	\begin{flushright}
		\begin{minipage}[t]{0.50\textwidth}
			%INFO JACKPOT CODING
			\begin{flushright}
				Gruppo: {\Large Jackpot Coding}\\
				Email: \href{mailto:jackpotcoding@gmail.com}{jackpotcoding@gmail.com}
			\end{flushright}
		\end{minipage}
	\end{flushright}
	%END JACKPOT CODING INFO
	
	\vspace{24pt}
	
	%TITLE
	\begin{center}
		\textbf{\LARGE SPECIFICA ARCHITETTURALE}
	\end{center}
	%END TITLE
	
	\vspace{13pt}
	
	\begin{flushleft}
		\begin{spacing}{1.5}
			REDATTORI: G. Moretto, M.Gobbo\\%INSERT HERE THE NAMES
			VERIFICATORI:\\
			\vspace{7pt}
			DESTINATARI: Prof. T. Vardanega, Prof. R. Cardin\\%INSERT HERE THE NAMES
		\end{spacing}
	\end{flushleft}
	
	\begin{flushright}
		\begin{spacing}{1}
			USO: ESTERNO\\
			VERSIONE: 0.0.4\\
		\end{spacing}
	\end{flushright}
	
	
	% Restore original margins from the second page onwards
	\restoregeometry
	
	\pagebreak
	
	\textbf{\Large Registro delle modifiche}
	\begin{table}[H]
		\centering
		\rowcolors{2}{black!15}{}
		\resizebox{\linewidth}{!}{
			\begin{tabular}{|c|c|c|c|c|}
				\rowcolor{teal!50}
				\hline
				\textbf{Versione} & \textbf{Data} & \textbf{Autore} & \textbf{Verificatore} & \textbf{Modifica} \\
    				\hline
				0.0.4 & 04/04/2024 & M. Camillo & - & Introduzione Documento \\
                		\hline
				0.0.4 & 04/04/2024 & M. Gobbo & - & Introduzione all'architettura \\
				\hline
				0.0.3 & 03/04/2024 & G. Moretto & M. Gobbo & Aggiunto elenco delle tabelle \\
				\hline
				0.0.2 & 02/04/2024 & G. Moretto & M. Gobbo & Aggiunta tabelle tecnologie codifica e testing \\
				\hline
				0.0.1 & 27/03/2024 & G. Moretto & M. Gobbo & Aggiunta struttura documento \\
				\hline
			\end{tabular}
		}
		\label{tab:conference}
	\end{table}
	
	\pagebreak
	\tableofcontents
	\pagebreak
	
	\section{Introduzione}
	
	\subsection{Scopo del Documento}

    Lo scopo di questo documento, \textit{Specifica Architetturale}, è quello di presentare, motivandole, le scelte architetturali e di design che è stato deciso di applicare al prodotto, oltre che a fornire una lista completa delle tecnologie utilizzate. La struttura interna del prodotto è esposta all'interno del dcoumento sotto forma di diagramma delle classe, in maniera da rendere più chiaro il software sviluppato. Vengono inoltre approfonditi e motivati a loro volta i design pattern utilizzati. 
	
	\subsection{Scopo del Prodotto}
    Il capitolato proposto dall'azienda \textit{Zucchetti S.p.A.} ha come obiettivo la realizzazione di un applicativo web al fine di ottenere un riscontro riguardo la fattibilità di un prodotto in grado di elaborare una frase in linguaggio naturale, fornita da un utente inesperto, e genere come output una \textit{query SQL} in grado di interrogare un database, di cui è conosciuta la struttura, in modo efficiente e che fornisca le informazioni richieste dall'utente.
	
	\subsection{Glossario}
    Al fine di evitare ambiguità o incomprensioni relative alla terminologia usata all'interno del documento, è fornito un \textit{Glossario} in cui vengono riportate definizioni precise per ogni termine potenzialmente ambiguo. La presenza di tali termini all'interno del documento è indicata con la presenza, vicino alla voce, di una \textit{G} in apice ($^G$). 
	\subsection{Riferimenti}
	
	\section{Tecnologie}
	
	\subsection{Codifica}

	\begin{longtblr}[
			caption = {Tecnologie di codifica.},
		]
		{
			colspec={|Q[0.15\linewidth]|Q[0.15\linewidth]|Q[0.70\linewidth]|},
			rows={halign=l},
			column{1}={halign=c},
			column{2}={halign=c},
			column{3}={halign=l},
			row{1}={halign=c},
			row{odd} = {gray!20},
			row{1}={teal!50},
			row{2}={teal!50},
			row{6}={teal!50},
			row{10}={teal!50}
		}
	
		\hline
		\textbf{Tecnologia} & \textbf{Versione} & \textbf{Descrizione} \\
		\hline
		\SetCell[c=3]{c} \textbf{Linguaggi} \\
		\hline
		HTML & 5 & Linguaggio di \textit{markup} utilizzato per la definizione della struttura di pagine \textit{web} \\
		\hline
		CSS & 3 & Linguaggio utilizzato per applicare stile a elementi presenti in una pagina \textit{HTML} \\
		\hline
		Python & 3.11.8 & Linguaggio di programmazione ad alto livello, orientato agli oggetti. Viene utilizzato per la creazione del \textit{server}. \\
		\hline
		\SetCell[c=3]{c} \textbf{Framework e Librerie} \\
		\hline
		Django & 5.0.3 & \textit{Framework} per la creazione di applicazioni \textit{web} scritto in linguaggio \textit{Python}. \\
		\hline
		TensorFlow & 2.15.0 & Libreria \textit{Python} per l'apprendimento automatico. \\
		\hline
		Transformers & 4.29.3 & Libreria \textit{Python} per l'utilizzo di modelli del portale \textit{Hugging Face} utilizzando \textit{TensorFlow}\\
		\hline
		\SetCell[c=3]{c} \textbf{Strumenti} \\
		\hline
		Pip & 24.0 & Strumento per la gestione dei pacchetti utilizzati da applicazioni \textit{Python}.\\
		\hline
		Git & 2.44.0 & Strumento per il controllo di versione utilizzato per la gestione della \textit{repository} remota presente su \textit{GitHub}. \\
		\hline
	\end{longtblr}
	
	\subsection{Testing}
	\begin{longtblr}[
		caption = {Tecnologie di testing.},
		]
		{
			colspec={|Q[0.15\linewidth]|Q[0.15\linewidth]|Q[0.70\linewidth]|},
			rows={halign=l},
			column{1}={halign=c},
			column{2}={halign=c},
			column{3}={halign=l},
			row{1}={halign=c},
			row{odd} = {gray!20},
			row{1}={teal!50},
			row{2}={teal!50},
			row{7}={teal!50}
		}
		\hline
		\textbf{Tecnologia} & \textbf{Versione} & \textbf{Descrizione} \\
		\hline
		\SetCell[c=3]{c} \textbf{Framework e Librerie} \\
		\hline
		Unittest & 3.11.8 & \textit{Framework} incluso nel linguaggio \textit{Python} utilizzato per il \textit{testing} di unità, utilizzato dal \textit{framework Django}.\\
		\hline
		Django Test Client & 5.0.3 & \textit{Client} per il \textit{testing} di un applicazione \textit{web} simulando un \textit{browser}, integrato nel \textit{framework Django}.\\
		\hline
		coverage.py & 7.4.4 & \textit{Tool} per misurare il \textit{code coverage} in applicazioni \textit{Python} integrabile nel \textit{framework Django}. \\
		\hline
		Prospector & 1.10.3 & \textit{Tool} per l'analisi statica di codice scritto nel linguaggio \textit{Python}. \\
		\hline
		\SetCell[c=3]{c} \textbf{Strumenti} \\
		\hline
		GitHub Actions & - & Servizio di \textit{Github} per la \textit{Continuous Integration}, automatizza i processi di \textit{build, test e deploy} del prodotto \textit{software}.\\
		\hline
	\end{longtblr}
	
	\section{Architettura}
	
	\subsection{Introduzione}
L'architettura di \textit{"Chat SQL"} si basa sul \textit{design pattern} architetturale \textit{MVT(Model View Template)}.\\
In questo modello, il \textit{"Model"} rappresenta i dati e la logica di \textit{business} dell'applicazione, la \textit{"View"} è responsabile della presentazione dei dati all'utente e la \textit{"Template"} definisce la struttura della \textit{View} e come i dati vengono visualizzati al suo interno. Più nello specifico il \textit{"Template"}, nel caso della nostra applicazione, definisce la struttura \textit{HTML}.\\
Questo \textit{design pattern} è simile al classico \textit{MVC(Model View Controller)} soltanto che al posto della parte di \textit{Controller} c'è il \textit{Template}.\\\\
E' stato scelta questa architettura in quanto offre i seguenti vantaggi:
\begin{itemize}
    \item \textbf{Separazione delle responsabilità}: \textit{MVT} separa chiaramente le responsabilità tra \textit{Model}, \textit{View} e \textit{Template}. Questo permette una migliore organizzazione del codice, facilitando la manutenzione e la scalabilità dell'applicazione.

    \item \textbf{Riutilizzo dei template}: I \textit{template} in \textit{MVT} consentono di separare la presentazione dalla logica di \textit{business}. Ciò facilita il riutilizzo dei componenti di interfaccia utente in diverse parti dell'applicazione, riducendo la duplicazione del codice e migliorando l'efficienza dello sviluppo.

    \item \textbf{Aumento delle Performance}: Utilizzare i \textit{template} pre-renderizzati può migliorare le prestazioni dell'applicazione rispetto a un'architettura \textit{MVC} tradizionale, poiché il \textit{rendering} dei \textit{template} può essere più efficiente rispetto ad esempio alla generazione dinamica di \textit{HTML} lato \textit{server}.

\end{itemize}
	
	\subsection{Diagrammi delle classi}
	
	\subsection{Design pattern utilizzati}
	
	\section{Elenco Figure}
	
	\section{Eleco tabelle}
	\makeatletter
	\@starttoc{lot}% Print List of Tables
	\makeatother
	
	
\end{document}
